% ================= §7.4 射影坐标和射影坐标变换 =================
\Needspace{12\baselineskip}
\section{射影坐标和射影坐标变换}

\S7.1 的齐次坐标借用了欧氏平面上的仿射标架 $[O_1;\vv d_1,\vv d_2]$。我们希望\key{直接用射影平面上的“点”做标架}。

\subsection{点的射影坐标}

先在把 $O$ 里想：取仿射标架 $[O;\vv d_1,\vv d_2,\vv d_3]$，把中每条直线就有了坐标（方向向量的坐标，差一个倍数）。但“点”只是直线，\key{方向向量的长度没有意义}：三条基本直线 $l_1,l_2,l_3$ 只确定 $\vv d_i$ 的方向，各自的倍数无法确定——这时坐标就不唯一了。

\key{补救办法}：再取一条直线 $l_4$，规定 $\vv d=\vv d_1+\vv d_2+\vv d_3$ 是 $l_4$ 的方向向量。若换取 $\vv d'=\lambda\vv d$，$\vv d_i'=\lambda_i\vv d_i$ 且 $\vv d'=\sum\vv d_i'$，比较系数得 $\lambda_1=\lambda_2=\lambda_3=\lambda$——各基向量的倍数被\key{统一}了，坐标只差一个公共倍数。

\begin{defbox}{定义 4.1–4.2（射影标架与射影坐标）}
$\pib$ 上取\key{一般位置的四个点}（任三点不共线）$A_1,A_2,A_3,E$。取各点对应的向量 $\vv a_1,\vv a_2,\vv a_3,\vv e$，并调整倍数使
\[
\vv e=\vv a_1+\vv a_2+\vv a_3 .
\]
若点 $M$ 对应 $\vv m=x_1\vv a_1+x_2\vv a_2+x_3\vv a_3$，则 $(x_1,x_2,x_3)\T$ 称为 $M$ 在\key{基底} $[A_1,A_2,A_3,E]$ 中的（齐次）\key{射影坐标}。$A_1,A_2,A_3$ 叫\key{基本点}，坐标 $(1,0,0)\T,(0,1,0)\T,(0,0,1)\T$；$E$ 叫\key{单位点}，坐标 $(1,1,1)\T$；$\triangle A_1A_2A_3$ 叫\key{坐标三角形}。
\end{defbox}

\begin{gdbox}{高代联动：为什么射影标架要 4 个点}
线性空间 $\R^3$ 的基是 3 个向量。但射影平面的“点”只是 1 维子空间，3 个点只给出 3 条“基向量所在的直线”，基向量可以各自伸缩——相当于基 $\{\vv a_1,\vv a_2,\vv a_3\}$ 只确定到 $\operatorname{diag}(k_1,k_2,k_3)$。第四个点 $E$（“单位点”）规定 $\vv e=\vv a_1+\vv a_2+\vv a_3$，把 $k_1=k_2=k_3$ 强制统一，剩下的公共倍数正是齐次坐标本来就允许的。一句话：\key{射影标架 = 基 + 比例尺}。
\end{gdbox}

\begin{figure}[htbp]
\centering
\begin{tikzpicture}[scale=1.15]
  \coordinate (A1) at (0,0); \coordinate (A2) at (6,0); \coordinate (A3) at (2.2,3.8);
  \coordinate (E) at (2.6,1.4); \coordinate (M) at (3.6,1.0);
  \coordinate (A12) at (intersection of A3--E and A1--A2);
  \coordinate (M12) at (intersection of A3--M and A1--A2);
  \coordinate (A23) at (intersection of A1--E and A2--A3);
  \coordinate (M23) at (intersection of A1--M and A2--A3);
  \fill[defbg] (A1)--(A2)--(A3)--cycle;
  \draw[thick,defblue] (A1)--(A2)--(A3)--cycle;
  \draw[figgray] (A3)--(A12) (A1)--(A23);
  \draw[figred,dashed] (A3)--(M12) (A1)--(M23);
  \foreach \p/\l/\pos in {A1/A_1/below left,A2/A_2/below right,A3/A_3/above}{\fill (\p) circle (1.6pt) node[\pos] {$\l$};}
  \fill (E) circle (1.6pt) node[above left] {$E$};
  \fill[figred] (M) circle (1.7pt) node[above right,figred] {$M$};
  \fill (A12) circle (1.4pt) node[below] {\scriptsize $A_{12}$};
  \fill[figred] (M12) circle (1.4pt) node[below,figred] {\scriptsize $M_{12}$};
  \fill (A23) circle (1.4pt) node[above right] {\scriptsize $A_{23}$};
  \fill[figred] (M23) circle (1.4pt) node[right,figred] {\scriptsize $M_{23}$};
  \node[align=left,right] at (6.4,1.9) {\small $x_1:x_2=(A_1,A_2;A_{12},M_{12})$\\[2pt]\small $x_2:x_3=(A_2,A_3;A_{23},M_{23})$\\[2pt]\small $x_3:x_1=(A_3,A_1;A_{31},M_{31})$};
\end{tikzpicture}
\caption{射影坐标的几何意义：坐标之比都是交比（$A_{12}=A_3E\cap A_1A_2$，$M_{12}=A_3M\cap A_1A_2$，……）}
\end{figure}

\begin{thmbox}{射影坐标的几何意义与两点注意}
\begin{itemize}
\item 对不在坐标三角形边上的点 $M$：$x_1:x_2=(A_1,A_2;A_{12},M_{12})$ 等（见上图），所以射影坐标\key{与 $O$ 的取法无关}，是射影平面自身的概念。
\item \S7.1 的齐次（仿射）坐标是\key{特殊基底}：$A_1,A_2$ 是 $x,y$ 轴的无穷远点，$A_3$ 是原点，$E$ 是点 $(1,1)$。
\item 一般基底下，\key{不能用 $x_3=0$ 判断无穷远点}——射影坐标根本“看不见”哪条线是无穷远直线。
\end{itemize}
\end{thmbox}

\subsection{射影坐标变换公式}

设两个基底 I$[A_1,A_2,A_3,E]$ 和 II$[A_1',A_2',A_3',E']$，点 $M$ 的 I 坐标 $X=(x_1,x_2,x_3)\T$，II 坐标 $X'$。则
\begin{equation}
\rho X=HX',\qquad \rho\neq0,\ |H|\neq0,
\end{equation}
其中 $H$ 的第 $j$ 列是 $\lambda_j\times(A_j'$ 的 I 坐标$)$，$\lambda_j$ 由
\[
\lambda\,(E'\text{ 的 I 坐标})=\lambda_1(A_1'\text{ 的 I 坐标})+\lambda_2(A_2'\text{ 的 I 坐标})+\lambda_3(A_3'\text{ 的 I 坐标})
\]
确定（与线性代数的“过渡矩阵的列 = 新基在旧基下的坐标”完全一样，只是新基向量的倍数要用 $E'$ 来校准）。

\begin{mtbox}{求射影坐标变换的三步法}
\textbf{①} 写出 $A_1',A_2',A_3',E'$ 的 I 坐标 $\vv a_1',\vv a_2',\vv a_3',\vv e'$；\quad
\textbf{②} 解 $\vv e'=\lambda_1\vv a_1'+\lambda_2\vv a_2'+\lambda_3\vv a_3'$（即解 $(\vv a_1',\vv a_2',\vv a_3')\vv\lambda=\vv e'$）；\\
\textbf{③} $H=(\lambda_1\vv a_1',\ \lambda_2\vv a_2',\ \lambda_3\vv a_3')$，得 $\rho X=HX'$。\\[2pt]
\textbf{检验}：$H(1,1,1)\T=\vv e'$；$\lambda_j$ 全不为零 $\iff$ 四点一般位置（某个 $\lambda_j=0$ 说明 $E'$ 落在坐标三角形的边上）。
\end{mtbox}

\begin{mtbox}{例 4.1}
$A_1',A_2',A_3',E'$ 的 I 坐标为 $(1,-1,2)\T,(2,0,1)\T,(-1,2,4)\T,(1,-1,0)\T$。解
\[
\lambda(1,-1,0)\T=\lambda_1(1,-1,2)\T+\lambda_2(2,0,1)\T+\lambda_3(-1,2,4)\T
\]
得 $\lambda_1=\frac7{15}\lambda,\ \lambda_2=\frac2{15}\lambda,\ \lambda_3=-\frac4{15}\lambda$，取 $\lambda=15$：$\lambda_1=7,\lambda_2=2,\lambda_3=-4$。于是
\[
\rho\begin{pmatrix}x_1\\x_2\\x_3\end{pmatrix}=\begin{pmatrix}7&4&4\\-7&0&-8\\14&2&-16\end{pmatrix}\begin{pmatrix}x_1'\\x_2'\\x_3'\end{pmatrix}.
\]
检验：$H(1,1,1)\T=(15,-15,0)\T\sim(1,-1,0)\T=E'$。\checkmark
\end{mtbox}

\subsection{直线的射影坐标方程、用射影坐标计算交比}
与 \S7.1 完全一样：直线 $PQ$ 的参数方程 $\vv x=\lambda\vv p+\mu\vv q$，方程 $\eta_1x_1+\eta_2x_2+\eta_3x_3=0$（$\vv\eta=\vv p\times\vv q$），三点共线 $\iff$ 行列式 $=0$。

\begin{thmbox}{定理 4.1（交比公式在任意射影坐标下成立）}
在任一基底中，若 $\vv c=\lambda_1\vv a+\mu_1\vv b$，$\vv d=\lambda_2\vv a+\mu_2\vv b$，则 $(A,B;C,D)=\dfrac{\lambda_2}{\mu_2}\Big/\dfrac{\lambda_1}{\mu_1}$。
\end{thmbox}
\noindent\textbf{理由}：换到一个使公式已知成立的基底 II（$B_1,B_2$ 为无穷远点，即齐次坐标），坐标变换 $\rho X=HX'$ 是线性的，$\vv c=\lambda_1\vv a+\mu_1\vv b$ 变成 $\rho_3^{-1}H\vv c'=\lambda_1\rho_1^{-1}H\vv a'+\mu_1\rho_2^{-1}H\vv b'$，系数比 $\lambda_1/\mu_1$ 只多出因子 $\rho_2/\rho_1$，对 $D$ 也多出同一个因子，相除抵消。

\subsection{点的非齐次射影坐标}
对不在 $A_1A_2$ 上的点，令 $x=\frac{x_1}{x_3},\ y=\frac{x_2}{x_3}$，称为\key{非齐次射影坐标}（它们就是交比：$x=(A_1,A_3;A_{13},M_{13})$，$y=(A_2,A_3;A_{23},M_{23})$）。两个基底间的变换是\key{分式线性函数}：
\begin{equation}
x=\frac{a_{11}x'+a_{12}y'+a_{13}}{a_{31}x'+a_{32}y'+a_{33}},\qquad
y=\frac{a_{21}x'+a_{22}y'+a_{23}}{a_{31}x'+a_{32}y'+a_{33}}.
\end{equation}

\begin{intbox}{直觉：分母就是“地平线”}
仿射坐标变换的分母恒为 1；射影坐标变换多了分母 $a_{31}x'+a_{32}y'+a_{33}$。分母为零的那条直线被送到“无穷远”——这就是照片里的地平线。手机“文档扫描”把斜拍的纸张拉正，用的正是这种分式线性变换（\S7.5）。
\end{intbox}

% ---------------- 习题 7.4 ----------------
\xiti{7.4}

\begin{ex}{7.4.1}
设在射影平面的一个基底 I 中给定了点 $B_1(4,-2,3)\T$，$B_2(5,2,0)\T$，$B_3(1,3,-2)\T$，$F(1,1,0)\T$，$M(1,1,-1)\T$，证明点 $B_1,B_2,B_3,F$ 是一般位置的四个点，并且求基底 I 到基底 II$[B_1,B_2,B_3,F]$ 的射影坐标变换公式，以及 $M$ 在基底 II 中的射影坐标。
\end{ex}
\jie \textbf{一般位置}：$|\vv b_1,\vv b_2,\vv b_3|=\begin{vmatrix}4&5&1\\-2&2&3\\3&0&-2\end{vmatrix}=3\neq0$，三基本点不共线。解 $\vv f=\lambda_1\vv b_1+\lambda_2\vv b_2+\lambda_3\vv b_3$：
\[
\begin{cases}4\lambda_1+5\lambda_2+\lambda_3=1\\-2\lambda_1+2\lambda_2+3\lambda_3=1\\3\lambda_1-2\lambda_3=0\end{cases}
\Rightarrow\ \lambda_1=2,\ \lambda_2=-2,\ \lambda_3=3 .
\]
三个 $\lambda$ 都不为零，$F$ 不在坐标三角形的边上，所以四点是一般位置。

\textbf{变换公式}：$H=(2\vv b_1,-2\vv b_2,3\vv b_3)$，
\[
\rho\begin{pmatrix}x_1\\x_2\\x_3\end{pmatrix}=\begin{pmatrix}8&-10&3\\-4&-4&9\\6&0&-6\end{pmatrix}\begin{pmatrix}x_1'\\x_2'\\x_3'\end{pmatrix}.
\]
\textbf{$M$ 的 II 坐标}：解 $HX'=(1,1,-1)\T$，得 $X'=-\frac16(7,8,6)\T$，即 $M$ 的 II 坐标为 $(7,8,6)\T$。检验：$H(7,8,6)\T=(-6,-6,6)\T\sim(1,1,-1)\T$。\checkmark

\begin{ex}{7.4.2}
在射影平面上，求基底 I$[A_1,A_2,A_3,E]$ 到基底 II$[A_3,A_1,E,A_2]$ 的射影坐标变换公式。
\end{ex}
\jie 新基底的四个点在 I 中的坐标：$A_1'=A_3=(0,0,1)\T$，$A_2'=A_1=(1,0,0)\T$，$A_3'=E=(1,1,1)\T$，$E'=A_2=(0,1,0)\T$。解
\[
(0,1,0)\T=\lambda_1(0,0,1)\T+\lambda_2(1,0,0)\T+\lambda_3(1,1,1)\T\ \Rightarrow\ \lambda_3=1,\ \lambda_2=-1,\ \lambda_1=-1 .
\]
所以
\[
\rho\begin{pmatrix}x_1\\x_2\\x_3\end{pmatrix}=\begin{pmatrix}0&-1&1\\0&0&1\\-1&0&1\end{pmatrix}\begin{pmatrix}x_1'\\x_2'\\x_3'\end{pmatrix}
\quad\Big(\text{与教材答案 }\begin{pmatrix}0&1&-1\\0&0&-1\\1&0&-1\end{pmatrix}\text{ 只差因子 }-1\Big).
\]
检验：$X'=(1,1,1)\T$（II 的单位点）$\mapsto X=(0,1,0)\T=A_2$。\checkmark

\begin{ex}{7.4.3}
在射影平面上，设点 $F$ 在基底 I$[A_1,A_2,A_3,E]$ 中的射影坐标为 $(\varepsilon_1,\varepsilon_2,\varepsilon_3)\T$，并且 $A_1,A_2,A_3,F$ 是一般位置的四个点，求基底 I 到基底 II$[A_1,A_2,A_3,F]$ 的射影坐标变换公式。
\end{ex}
\jie 基本点不变，只换单位点。$\vv f=\varepsilon_1\vv a_1+\varepsilon_2\vv a_2+\varepsilon_3\vv a_3$，所以 $\lambda_j=\varepsilon_j$：
\[
\rho X=\operatorname{diag}(\varepsilon_1,\varepsilon_2,\varepsilon_3)\,X' .
\]
一般位置保证 $\varepsilon_1\varepsilon_2\varepsilon_3\neq0$。\key{直觉}：换单位点 = 三个基向量各自伸缩——对角矩阵。

\begin{ex}{7.4.4}
在射影平面上取四点 $B_1(1,2,1)\T,B_2(1,1,0)\T,B_3(2,1,1)\T,F(0,1,7)\T$，求点 $M(1,1,1)\T$ 在基底 $[B_1,B_2,B_3,F]$ 中的坐标。
\end{ex}
\jie $|\vv b_1,\vv b_2,\vv b_3|=-2\neq0$。解 $\vv f=\lambda_1\vv b_1+\lambda_2\vv b_2+\lambda_3\vv b_3$ 得 $\lambda_1=4,\lambda_2=-10,\lambda_3=3$（验：$4(1,2,1)-10(1,1,0)+3(2,1,1)=(0,1,7)$）。于是
\[
\rho X=\begin{pmatrix}4&-10&6\\8&-10&3\\4&0&3\end{pmatrix}X' .
\]
解 $HX'=(1,1,1)\T$ 得 $X'=(\frac18,\frac1{20},\frac16)\T\sim(15,6,20)\T$。检验：$H(15,6,20)\T=(120,120,120)\T$。\checkmark

\begin{ex}{7.4.5}
证明 Pappus（帕普斯）定理：设 $A,B,C$；$A',B',C'$ 分别是直线 $l_1$ 和 $l_2$ 上的三个点，并且均与 $l_1$ 和 $l_2$ 的交点 $D$ 不重合，又设 $AB'$ 与 $A'B$ 交于点 $P$，$BC'$ 与 $B'C$ 交于点 $Q$，$CA'$ 与 $C'A$ 交于点 $R$，则 $P,Q,R$ 三点共线。
\end{ex}

\begin{figure}[htbp]
\centering
\begin{tikzpicture}[scale=1.05]
  \coordinate (A) at (0.6,0); \coordinate (B) at (3.0,0); \coordinate (C) at (6.2,0);
  \coordinate (A') at (0.9,2.9); \coordinate (B') at (3.9,3.1); \coordinate (C') at (5.8,3.24);
  \draw[thick] (-0.3,0)--(7,0) node[right] {$l_1$};
  \draw[thick] (-0.3,2.82)--(7,3.32) node[right] {$l_2$};
  \draw[figblue] (A)--(B') (B)--(A');
  \draw[figred] (B)--(C') (C)--(B');
  \draw[figgreen!70!black] (C)--(A') (A)--(C');
  \coordinate (P) at (intersection of A--B' and A'--B);
  \coordinate (Q) at (intersection of B--C' and B'--C);
  \coordinate (R) at (intersection of C--A' and C'--A);
  \draw[keyred,very thick] ($(P)!-0.15!(Q)$)--($(P)!1.25!(Q)$);
  \foreach \p/\pos in {A/below,B/below,C/below,A'/above,B'/above,C'/above}{\fill (\p) circle (1.6pt) node[\pos] {$\p$};}
  \foreach \p in {P,Q,R}{\fill[keyred] (\p) circle (2pt) node[below=2pt,keyred] {$\p$};}
  \node[keyred] at ($(P)!1.36!(Q)+(0,0.1)$) {\small Pappus 线};
\end{tikzpicture}
\caption{Pappus 定理：两条直线上各三点“交叉连线”的三个交点共线}
\end{figure}

\zheng 取基底 $[O_1,O_2,D,E]$，其中 $O_1\in l_1$、$O_2\in l_2$ 与六个点都不同，$E$ 为一般位置的点。于是 $O_1=(1,0,0)\T$，$O_2=(0,1,0)\T$，$D=(0,0,1)\T$，$l_1=O_1D:\ x_2=0$，$l_2=O_2D:\ x_1=0$。$l_1$ 上异于 $D$ 的点可写成 $(1,0,t)\T$，$l_2$ 上的写成 $(0,1,t)\T$：
\[
A(1,0,a),\ B(1,0,b),\ C(1,0,c);\qquad A'(0,1,a'),\ B'(0,1,b'),\ C'(0,1,c').
\]
用叉积：$AB'=(-a,-b',1)\T$，$A'B=(b,a',-1)\T$，
\[
\vv p=AB'\times A'B=(b'-a',\ b-a,\ bb'-aa')\T .
\]
把 $(A,B)$ 依次换成 $(B,C)$、$(C,A)$ 得
\[
\vv q=(c'-b',\ c-b,\ cc'-bb')\T,\qquad \vv r=(a'-c',\ a-c,\ aa'-cc')\T .
\]
显然 $\vv p+\vv q+\vv r=\vv0$（每个分量都是“循环差”之和），三向量线性相关，$P,Q,R$ 共线。\qed

\begin{intbox}{直觉}
和 Desargues 的证明一样，结论又是“三个向量之和为零”。选 $D$ 为基本点、两条直线为坐标三角形的两边，是把“两条直线”的条件吃进坐标里的标准手法。（把 $l_1\cup l_2$ 看成一条退化的二次曲线，Pappus 就是 \S7.6 Pascal 定理的退化情形。）
\end{intbox}

\begin{ex}{7.4.6}
在射影平面上取一个基底 $[A_1,A_2,A_3,E]$，设 $A_{23}$ 是直线 $A_1E$ 与 $A_2A_3$ 的交点，求直线 $A_2A_3$ 上的一个点 $P$，使得交比 $(A_2,A_3;A_{23},P)=a$。
\end{ex}
\jie $A_1E:\ x_2=x_3$，$A_2A_3:\ x_1=0$，所以 $A_{23}=(0,1,1)\T=\vv a_2+\vv a_3$（$\frac{\lambda_1}{\mu_1}=1$）。设 $P=\lambda\vv a_2+\mu\vv a_3$，由定理 4.1，$\frac{\lambda}{\mu}\big/1=a$，取 $\lambda=a,\mu=1$：
\[
P=(0,a,1)\T .
\]
这也正是“非齐次射影坐标就是交比”的一个具体例子：$A_2A_3$ 上点 $(0,x_2,x_3)$ 的 $x_2:x_3$ 就等于 $(A_2,A_3;A_{23},P)$。

\begin{ex}{7.4.7}
在射影平面上给了一个三角形 $A_1A_2A_3$，设 $P_1,P_2,P_3$ 分别在直线 $A_2A_3,A_3A_1,A_1A_2$ 上，但是都不和 $A_1,A_2,A_3$ 重合，而 $Q_1,Q_2,Q_3$ 依次为 $P_1,P_2,P_3$ 对于 $A_2,A_3$；$A_3,A_1$；$A_1,A_2$ 的调和共轭点，证明：$A_1Q_1,A_2Q_2,A_3Q_3$ 共点的充分必要条件是 $P_1,P_2,P_3$ 共线。
\end{ex}

\begin{figure}[htbp]
\centering
\begin{tikzpicture}[scale=0.95]
  \coordinate (A1) at (0,0); \coordinate (A2) at (5.91,0); \coordinate (A3) at (4.0,2.54);
  \coordinate (O) at (4.79,0.6);
  \coordinate (Q1) at (5.41,0.67); \coordinate (Q2) at (2.69,1.71); \coordinate (Q3) at (5.03,0);
  \coordinate (P1) at (6.99,-1.43); \coordinate (P2) at (7.8,4.96); \coordinate (P3) at (7.17,0);
  \fill[defbg] (A1)--(A2)--(A3)--cycle;
  \draw[thick] (A1)--(P3) (A3)--(P1) (A1)--(P2);
  \draw[defblue,thick] (A1)--(A2)--(A3)--cycle;
  \draw[figblue] (A1)--(Q1) (A2)--(Q2) (A3)--(Q3);
  \draw[figred,thick] ($(P1)!-0.1!(P2)$)--($(P1)!1.08!(P2)$);
  \foreach \p/\l/\pos in {A1/A_1/below left,A2/A_2/below,A3/A_3/above left}{\fill (\p) circle (1.6pt) node[\pos] {$\l$};}
  \foreach \p/\l/\pos in {Q1/Q_1/right,Q2/Q_2/above left,Q3/Q_3/below}{\fill[figblue] (\p) circle (1.6pt) node[\pos,figblue] {\small $\l$};}
  \foreach \p/\l/\pos in {P1/P_1/right,P2/P_2/right,P3/P_3/below right}{\fill[figred] (\p) circle (1.6pt) node[\pos,figred] {\small $\l$};}
  \fill[keyred] (O) circle (1.8pt);
\end{tikzpicture}
\caption{$P_i$ 共线（红） $\iff$ 它们的调和共轭点 $Q_i$ 与对顶点的连线共点（蓝）}
\end{figure}

\zheng 取基底 $[A_1,A_2,A_3,E]$（$E$ 任取一般位置点）。设
\[
P_1=\vv a_2+t_1\vv a_3=(0,1,t_1)\T,\quad P_2=\vv a_3+t_2\vv a_1=(t_2,0,1)\T,\quad P_3=\vv a_1+t_3\vv a_2=(1,t_3,0)\T,
\]
$t_i\neq0$（不与顶点重合）。$P_1=1\cdot\vv a_2+t_1\vv a_3$，其调和共轭点为 $Q_1=\vv a_2-t_1\vv a_3$（系数比变号，交比 $=-1$）。同理
\[
Q_1=(0,1,-t_1)\T,\qquad Q_2=(-t_2,0,1)\T,\qquad Q_3=(1,-t_3,0)\T .
\]
\textbf{$P_i$ 共线}：$\begin{vmatrix}0&1&t_1\\t_2&0&1\\1&t_3&0\end{vmatrix}=1+t_1t_2t_3=0$。

\textbf{$A_iQ_i$ 共点}：三条线坐标 $A_1Q_1=(0,t_1,1)\T$，$A_2Q_2=(1,0,t_2)\T$，$A_3Q_3=(t_3,1,0)\T$，
\[
\begin{vmatrix}0&t_1&1\\1&0&t_2\\t_3&1&0\end{vmatrix}=1+t_1t_2t_3=0 .
\]
两个条件都是 $t_1t_2t_3=-1$，所以等价。\qed

\begin{intbox}{几何意义}
若 $P_1P_2P_3$ 是无穷远直线，则 $Q_i$ 是\key{各边中点}（中点 $\leftrightarrow$ 无穷远点调和共轭），结论变成“\key{三角形三条中线共点}”。一般情形是它的射影像：任意一条直线都可以被射影变换送到无穷远。
\end{intbox}

\begin{ex}{7.4.8}
在射影平面上，设 $A_1,A_2,A_3$ 和 $P_1,P_2,P_3$ 的意义同第 7 题，但是 $(A_2,A_3;Q_1,P_1)=e_1,\ (A_3,A_1;Q_2,P_2)=e_2,\ (A_1,A_2;Q_3,P_3)=e_3$，证明：若 $P_1,P_2,P_3$ 共线，则 $Q_1,Q_2,Q_3$ 共线的充分必要条件是 $e_1e_2e_3=1$。又如果 $P_1,P_2,P_3$ 在扩大的欧氏平面的无穷远直线上，那么上述命题的几何意义是什么？
\end{ex}
\zheng 沿用第 7 题的坐标，$P_1=\vv a_2+t_1\vv a_3$。设 $Q_1=\vv a_2+s_1\vv a_3$，则
\[
e_1=(A_2,A_3;Q_1,P_1)=\frac{1/t_1}{1/s_1}=\frac{s_1}{t_1},\quad\text{即}\quad s_1=e_1t_1;\qquad\text{同理 } s_2=e_2t_2,\ s_3=e_3t_3 .
\]
$Q_1=(0,1,s_1)\T,Q_2=(s_2,0,1)\T,Q_3=(1,s_3,0)\T$ 与 $P_i$ 同形，所以 $Q_i$ 共线 $\iff 1+s_1s_2s_3=0$。已知 $P_i$ 共线，$t_1t_2t_3=-1$，于是
\[
1+s_1s_2s_3=1+e_1e_2e_3\,t_1t_2t_3=1-e_1e_2e_3=0\iff e_1e_2e_3=1.\qquad\blacksquare
\]
\textbf{几何意义}：$P_i$ 在无穷远直线上时，$(A_2,A_3,P_1)=-1$，
\[
e_1=\frac{(A_2,A_3,Q_1)}{(A_2,A_3,P_1)}=-\frac{A_2Q_1}{Q_1A_3},\ \text{等},
\qquad e_1e_2e_3=1\iff \frac{A_2Q_1}{Q_1A_3}\cdot\frac{A_3Q_2}{Q_2A_1}\cdot\frac{A_1Q_3}{Q_3A_2}=-1,
\]
这就是 \key{Menelaus（梅涅劳斯）定理}（第一章 \S2）：三边（所在直线）上的三点共线的充要条件。

\begin{ex}{7.4.9}
在射影平面上，设 $A_1,A_2,A_3,P_1,P_2,P_3,Q_1,Q_2,Q_3$ 的意义同第 8 题，证明：若 $P_1,P_2,P_3$ 共线，则 $A_1Q_1,A_2Q_2,A_3Q_3$ 共点的充分必要条件是 $e_1e_2e_3=-1$。又如果 $P_1,P_2,P_3$ 在扩大的欧氏平面的无穷远直线上，那么上述命题的几何意义是什么？
\end{ex}
\zheng 线坐标 $A_1Q_1=\vv a_1\times\vv q_1=(0,-s_1,1)\T$，$A_2Q_2=(1,0,-s_2)\T$，$A_3Q_3=(-s_3,1,0)\T$，
\[
\begin{vmatrix}0&-s_1&1\\1&0&-s_2\\-s_3&1&0\end{vmatrix}=1-s_1s_2s_3 .
\]
共点 $\iff s_1s_2s_3=1\iff e_1e_2e_3\,t_1t_2t_3=1\iff -e_1e_2e_3=1\iff e_1e_2e_3=-1$。\qed

（第 7 题是 $e_1=e_2=e_3=-1$ 的特例，乘积 $=-1$，吻合。）

\textbf{几何意义}：$P_i$ 在无穷远直线上时，$e_1e_2e_3=-1\iff\frac{A_2Q_1}{Q_1A_3}\cdot\frac{A_3Q_2}{Q_2A_1}\cdot\frac{A_1Q_3}{Q_3A_2}=1$，这是 \key{Ceva（塞瓦）定理}（习题 2.3 第 12 题）。

\begin{ckbox}{对照：Menelaus 与 Ceva 是一对“射影孪生”}
同一个三角形、同样的分比 $\frac{A_2Q_1}{Q_1A_3}$，乘积 $=-1$ 是三点共线，$=+1$ 是三线共点。射影观点下差别仅在 $e_1e_2e_3=\pm1$。
\end{ckbox}

\begin{ex}{7.4.10$^*$}
射影平面 $\pib$ 上，由一般位置的四个点 $A,B,C,D$（称它们为顶点）和由它们两两相连的六条直线（称它们为边）构成的图形称为\key{完备四点形}，其中不在同一顶点上的两条边称为对边，对边的交点称为\key{对角点}，如图中的 $E,F,G$ 均是对角点。(1) 证明：完备四点形的三个对角点不共线。由这三个对角点确定的三角形称为完备四点形的\key{对角三角形}。(2) 证明：完备四点形的任意两个对角点调和分割它们的连线和另外两条边的交点。
\end{ex}

\begin{figure}[htbp]
\centering
\begin{tikzpicture}[scale=1.25]
  \coordinate (A) at (2.67,1.66); \coordinate (B) at (0.49,2.88); \coordinate (C) at (1.47,0.76); \coordinate (D) at (2.79,3.47);
  \coordinate (E) at (intersection of A--B and C--D);
  \coordinate (F) at (intersection of A--C and B--D);
  \coordinate (G) at (intersection of A--D and B--C);
  \coordinate (U) at (intersection of E--F and A--D);
  \coordinate (V) at (intersection of E--F and B--C);
  \draw[thick] ($(B)!-0.12!(A)$)--($(B)!1.15!(A)$);   % AB
  \draw[thick] ($(C)!-0.1!(D)$)--($(C)!1.1!(D)$);     % CD
  \draw[thick] ($(C)!-0.06!(F)$)--($(C)!1.06!(F)$);   % AC
  \draw[thick] ($(B)!-0.06!(F)$)--($(B)!1.06!(F)$);   % BD
  \draw[thick] ($(G)!-0.06!(D)$)--($(G)!1.08!(D)$);   % AD
  \draw[thick] ($(B)!-0.06!(G)$)--($(B)!1.08!(G)$);   % BC
  \draw[figred,very thick] ($(V)!-0.12!(F)$)--($(V)!1.06!(F)$);
  \draw[figgray,dashed] (E)--(G)--(F);
  \foreach \p/\pos in {A/below right,B/right,C/left,D/above}{\fill (\p) circle (1.7pt) node[\pos] {$\p$};}
  \foreach \p/\pos in {E/left,F/above right,G/below}{\fill[figblue] (\p) circle (2pt) node[\pos,figblue] {$\p$};}
  \fill[figred] (U) circle (1.7pt) node[above left=-1pt,figred] {\small $U$};
  \fill[figred] (V) circle (1.7pt) node[below=2pt,figred] {\small $V$};
\end{tikzpicture}
\caption{完备四点形 $ABCD$，对角点 $E,F,G$（蓝），$(E,F;U,V)=-1$}
\end{figure}

\zheng 取基底 $[A,B,C,D]$：$A=(1,0,0)\T,B=(0,1,0)\T,C=(0,0,1)\T,D=(1,1,1)\T$。六条边：
\[
AB:x_3=0,\ CD:x_1=x_2;\quad AC:x_2=0,\ BD:x_1=x_3;\quad AD:x_2=x_3,\ BC:x_1=0 .
\]
对角点：$E=AB\cap CD=(1,1,0)\T$，$F=AC\cap BD=(1,0,1)\T$，$G=AD\cap BC=(0,1,1)\T$。

(1) $\begin{vmatrix}1&1&0\\1&0&1\\0&1&1\end{vmatrix}=-2\neq0$，三个对角点不共线。

(2) 以 $E,F$ 为例（另两种情形同理，或换基底的顶点顺序）。$EF=(1,1,0)\T\times(1,0,1)\T=(1,-1,-1)\T$：$x_1-x_2-x_3=0$。另外两条边是不过 $E,F$ 的 $AD,BC$：
\[
U=EF\cap AD=(2,1,1)\T=\vv e+\vv f,\qquad V=EF\cap BC=(0,1,-1)\T=\vv e-\vv f,
\]
所以 $(E,F;U,V)=\dfrac{1/(-1)}{1/1}=-1$。\qed

\begin{intbox}{有趣的反例：Fano 平面}
(1) 的行列式等于 $-2$。在特征为 2 的域上 $-2=0$——$PG(2,2)$（Fano 平面）中完备四点形的对角点\key{共线}（\S7.1 图中那个“圆”）！所以“对角点不共线”其实依赖于 $1+1\neq0$。
\end{intbox}

\begin{ex}{7.4.11$^*$}
在射影平面上，完备四点形的对偶称为\key{完备四边形}。叙述第 10 题中第 (2) 小题的对偶命题，并且说明调和线束的作图方法。
\end{ex}
\jie\textbf{完备四边形}：一般位置（任三条不共点）的四条直线（边），它们两两的六个交点（顶点）；不在同一条边上的两个顶点叫\key{相对顶点}，三对相对顶点的连线叫\key{对角线}。

\begin{thmbox}{第 10(2) 题的对偶命题}
完备四边形的任意两条对角线，被它们的交点与另外两个顶点（即第三条对角线上的两个顶点）的连线调和分割。也就是：若对角线 $x,y$ 交于 $O$，第三条对角线上的顶点为 $W_1,W_2$，则 $(x,y;OW_1,OW_2)=-1$。
\end{thmbox}

\textbf{调和线束的作图}：已知共点于 $O$ 的三条直线 $l_1,l_2,l_3$，求第四调和线 $l_4$。
\begin{enumerate}
\item 作一条不过 $O$ 的直线 $l$，交 $l_i$ 于 $B_i$（$i=1,2,3$）；
\item 过 $B_3$ 作另一条直线 $l'$，交 $l_1,l_2$ 于 $D_1,D_2$；
\item 连 $D_1B_2$、$D_2B_1$，交于 $D_4$；
\item 直线 $OD_4$ 即为所求 $l_4$。
\end{enumerate}

\begin{figure}[htbp]
\centering
\begin{tikzpicture}[scale=1.05]
  \coordinate (O) at (0,4.2);
  \coordinate (E1) at ($(O)+(-112:5)$); \coordinate (E2) at ($(O)+(-62:5)$); \coordinate (E3) at ($(O)+(-42:6.6)$);
  \coordinate (La) at (-2.6,0.2); \coordinate (Lb) at (5.4,0.2);
  \coordinate (B1) at (intersection of O--E1 and La--Lb);
  \coordinate (B2) at (intersection of O--E2 and La--Lb);
  \coordinate (B3) at (intersection of O--E3 and La--Lb);
  \coordinate (K) at ($(B3)+(165:1)$);
  \coordinate (D1) at (intersection of O--E1 and B3--K);
  \coordinate (D2) at (intersection of O--E2 and B3--K);
  \coordinate (D4) at (intersection of D1--B2 and D2--B1);
  \draw[thick] (O)--(E1) node[below] {$l_1$}; \draw[thick] (O)--(E2) node[below] {$l_2$}; \draw[thick] (O)--(E3) node[below] {$l_3$};
  \draw[figgray] (La)--(Lb) node[right] {$l$};
  \draw[figgray] ($(B3)!-0.1!(D1)$)--($(B3)!1.12!(D1)$) node[left] {$l'$};
  \draw[figblue] (D1)--(B2) (D2)--(B1);
  \draw[figred,very thick] (O)--($(O)!1.25!(D4)$) node[below,figred] {$l_4$};
  \foreach \p/\l/\pos in {B1/B_1/below left,B2/B_2/below right,B3/B_3/below,D1/D_1/left,D2/D_2/above right}{\fill (\p) circle (1.6pt) node[\pos] {\small $\l$};}
  \fill (O) circle (1.7pt) node[above] {$O$};
  \fill[figred] (D4) circle (1.8pt) node[right=2pt,figred] {\small $D_4$};
\end{tikzpicture}
\caption{作第四调和线：$(l_1,l_2;l_3,l_4)=-1$}
\end{figure}

\zheng 看完备四点形 $B_1D_1B_2D_2$（顶点 $B_1,D_1\in l_1$，$B_2,D_2\in l_2$）。它的三个对角点是
\[
B_1D_1\cap B_2D_2=l_1\cap l_2=O,\qquad B_1B_2\cap D_1D_2=l\cap l'=B_3,\qquad B_1D_2\cap B_2D_1=D_4 .
\]
对对角点 $B_3,D_4$ 用第 10(2) 题：直线 $B_3D_4$ 与另外两条边 $B_1D_1=l_1$、$B_2D_2=l_2$ 交于 $Z_1,Z_2$，则 $(B_3,D_4;Z_1,Z_2)=-1$。从 $O$ 投影（$OB_3=l_3$，$OD_4=l_4$，$OZ_1=l_1$，$OZ_2=l_2$）：
\[
(l_3,l_4;l_1,l_2)=-1\ \Rightarrow\ (l_1,l_2;l_3,l_4)=-1 .\qquad\blacksquare
\]
（这正是 \S7.3 “作第四调和点”的对偶作图。）
